% GENERATED FILE. Edit canonical lessons, then run npm run generate:books.
% canonical-source-hash: fnv1a64:2a0c0f484b4274df
% canonical-lessons: IT-C287-diretta, IT-C287-trasmissione, IT-C287-trasmettere, IT-C287-canale, IT-C287-emittente

\chapter{In onda}
\label{ch:it-in-onda}
\hlchaptermodality{287}

\begin{chapteropening}
\textbf{By the end of this chapter:} \emph{I can talk about broadcasting: going out live, a programme, a channel and a station.}

Complete the last lesson of chapter 287: I can talk about broadcasting: going out live, a programme, a channel and a station.
\end{chapteropening}

\section[la diretta]{la diretta — the live broadcast}

\label{lesson:IT-C287-diretta}



\begin{quote}
Before the new one: say the Italian for censorship, then the Italian for to censor.
\end{quote}

\subsection*{You'll want to know: la diretta}
\textbf{la diretta} — \textquotedblleft{}the live broadcast\textquotedblright{}.

\begin{quote}
Guardiamo la partita in diretta.
\end{quote}

\begin{quote}
\emph{Let us watch the match live.}
\end{quote}

\emph{In diretta} — live; the opposite is \emph{in differita}, recorded.

\begin{cousinweb}
\emph{diretta} is the feminine of \emph{diretto} \textquotedblleft{}direct\textquotedblright{}, from Latin \emph{dīrēctus} \textquotedblleft{}straight\textquotedblright{}, the past participle of \emph{dīrigere} \textquotedblleft{}to set straight\textquotedblright{} (Lewis \& Short).
\end{cousinweb}

\begin{grammarlens}[title={What you've built}]
One more word for telling, explaining and describing at B1.
\end{grammarlens}

\subsection*{Your turn}
\begin{itemize}
\raggedright
  \item \emph{Say it:} \emph{la diretta}
  \item \emph{Say it:} \emph{la diretta}, once more
  \item \emph{Say it:} a sentence of your own with \emph{la diretta}
  \item \emph{Recall:} say the Italian for censorship, then the Italian for to censor, then say \emph{la diretta} again
\end{itemize}

\begin{tcolorbox}[breakable,colback=teal!4,colframe=teal!35!black,title={Before you move on}]
What does la diretta mean? (\textquotedblleft{}The live broadcast\textquotedblright{}.) Say it once more.
\end{tcolorbox}
\section[la trasmissione]{la trasmissione — the programme, the broadcast; transmission}

\label{lesson:IT-C287-trasmissione}



\begin{quote}
Before the new one: say the Italian for to censor, then the Italian for the live broadcast.
\end{quote}

\subsection*{You'll want to know: la trasmissione}
\textbf{la trasmissione} — \textquotedblleft{}the programme, the broadcast; transmission\textquotedblright{}.

\begin{quote}
È una trasmissione molto seguita.
\end{quote}

\begin{quote}
\emph{It is a very popular programme.}
\end{quote}

On TV or radio: \emph{una trasmissione di cucina}.

\begin{cousinweb}
\emph{trasmissione} is a learned borrowing from Latin \emph{trānsmissiō} \textquotedblleft{}a sending across\textquotedblright{} (EtymDB), from \emph{trānsmittere}.
\end{cousinweb}

\begin{grammarlens}[title={What you've built}]
One more word for telling, explaining and describing at B1.
\end{grammarlens}

\subsection*{Your turn}
\begin{itemize}
\raggedright
  \item \emph{Say it:} \emph{la trasmissione}
  \item \emph{Say it:} \emph{la trasmissione}, once more
  \item \emph{Say it:} a sentence of your own with \emph{la trasmissione}
  \item \emph{Recall:} say the Italian for to censor, then the Italian for the live broadcast, then say \emph{la trasmissione} again
\end{itemize}

\begin{tcolorbox}[breakable,colback=teal!4,colframe=teal!35!black,title={Before you move on}]
What does la trasmissione mean? (\textquotedblleft{}The programme, the broadcast; transmission\textquotedblright{}.) Say it once more.
\end{tcolorbox}
\section[trasmettere]{trasmettere — to broadcast; to pass on}

\label{lesson:IT-C287-trasmettere}



\begin{quote}
Before the new one: say the Italian for the live broadcast, then the Italian for the programme.
\end{quote}

\subsection*{You'll want to know: trasmettere}
\textbf{trasmettere} — \textquotedblleft{}to broadcast; to pass on\textquotedblright{}.

\begin{quote}
La radio trasmette musica tutta la notte.
\end{quote}

\begin{quote}
\emph{The radio broadcasts music all night.}
\end{quote}

Conjugated like \emph{mettere}: \emph{ha trasmesso}.

\begin{cousinweb}
\emph{trasmettere} is a learned borrowing from Latin \emph{trānsmittere} \textquotedblleft{}to send across\textquotedblright{} (EtymDB), from \emph{trāns} and \emph{mittere} \textquotedblleft{}to send\textquotedblright{}.
\end{cousinweb}

\begin{grammarlens}[title={What you've built}]
One more word for telling, explaining and describing at B1.
\end{grammarlens}

\subsection*{Your turn}
\begin{itemize}
\raggedright
  \item \emph{Say it:} \emph{trasmettere}
  \item \emph{Say it:} \emph{trasmettere}, once more
  \item \emph{Say it:} a sentence of your own with \emph{trasmettere}
  \item \emph{Recall:} say the Italian for the live broadcast, then the Italian for the programme, then say \emph{trasmettere} again
\end{itemize}

\begin{tcolorbox}[breakable,colback=teal!4,colframe=teal!35!black,title={Before you move on}]
What does trasmettere mean? (\textquotedblleft{}To broadcast; to pass on\textquotedblright{}.) Say it once more.
\end{tcolorbox}
\section[il canale]{il canale — the TV channel; the canal}

\label{lesson:IT-C287-canale}



\begin{quote}
Before the new one: say the Italian for the programme, then the Italian for to broadcast.
\end{quote}

\subsection*{You'll want to know: il canale}
\textbf{il canale} — \textquotedblleft{}the TV channel; the canal\textquotedblright{}.

\begin{quote}
Cambia canale, per favore!
\end{quote}

\begin{quote}
\emph{Change the channel, please!}
\end{quote}

In a city built on water, \emph{il canale} is a canal.

\begin{cousinweb}
\emph{canale} is inherited from Latin \emph{canālis} \textquotedblleft{}a pipe, a channel\textquotedblright{} (EtymDB), built on \emph{canna} \textquotedblleft{}reed\textquotedblright{}.
\end{cousinweb}

\begin{grammarlens}[title={What you've built}]
One more word for telling, explaining and describing at B1.
\end{grammarlens}

\subsection*{Your turn}
\begin{itemize}
\raggedright
  \item \emph{Say it:} \emph{il canale}
  \item \emph{Say it:} \emph{il canale}, once more
  \item \emph{Say it:} a sentence of your own with \emph{il canale}
  \item \emph{Recall:} say the Italian for the programme, then the Italian for to broadcast, then say \emph{il canale} again
\end{itemize}

\begin{tcolorbox}[breakable,colback=teal!4,colframe=teal!35!black,title={Before you move on}]
What does il canale mean? (\textquotedblleft{}The TV channel; the canal\textquotedblright{}.) Say it once more.
\end{tcolorbox}
\section[l'emittente]{l'emittente — the broadcaster, the station}

\label{lesson:IT-C287-emittente}



\begin{quote}
Before the new one: say the Italian for to broadcast, then the Italian for the TV channel.
\end{quote}

\subsection*{You'll want to know: l'emittente}
\textbf{l'emittente} — \textquotedblleft{}the broadcaster, the station\textquotedblright{}.

\begin{quote}
È un'emittente locale: trasmette solo in questa regione.
\end{quote}

\begin{quote}
\emph{It is a local station: it only broadcasts in this region.}
\end{quote}

Feminine: \emph{un'emittente}.

\begin{cousinweb}
\emph{emittente} is a learned word from Latin \emph{ēmittēns}, the present participle of \emph{ēmittere} \textquotedblleft{}to send out\textquotedblright{} (Lewis \& Short).
\end{cousinweb}

\begin{grammarlens}[title={What you've built}]
One more word for telling, explaining and describing at B1.
\end{grammarlens}

\subsection*{Your turn}
\begin{itemize}
\raggedright
  \item \emph{Say it:} \emph{l'emittente}
  \item \emph{Say it:} \emph{l'emittente}, once more
  \item \emph{Say it:} a sentence of your own with \emph{l'emittente}
  \item \emph{Recall:} say the Italian for to broadcast, then the Italian for the TV channel, then say \emph{l'emittente} again
\end{itemize}

\begin{tcolorbox}[breakable,colback=teal!4,colframe=teal!35!black,title={Before you move on}]
What does l'emittente mean? (\textquotedblleft{}The broadcaster, the station\textquotedblright{}.) Say it once more.
\end{tcolorbox}
